\chapter{N-Body Dynamics for Semantic Galaxies}
\label{ch:nbody}

The preceding chapters developed the geometric, sheaf-theoretic,
and embedding-based architecture for semantic galaxy morphogenesis.
We now extend this system by introducing an \emph{N-body formulation}
of Polyxan atoms evolving under RSVP field forces,
in analogy with gravitational or electrostatic simulations,
but enriched by curvature, density, and semantic constraints.

The N-body formulation serves two purposes:

\begin{enumerate}
\item It provides a scalable approximation to full embedding PDEs
      by modeling atom dynamics as particle-like flows.
\item It reveals emergent structures such as clusters, filaments,
      basins of attraction, and bifurcation events
      in semantic meaning-space.
\end{enumerate}

This chapter formalizes the particle equations of motion,
derives the coupled RSVP--Polyxan–sheaf forces,
proves stability-theoretic properties,
and provides discretization rules enabling GPU-accelerated simulation.


% ============================================================
\section{From Embedding Flow to N-Body Approximation}
% ============================================================

Recall from Chapter~\ref{ch:embedding-flow} that
\[
  \dot\iota(v)
    = -\frac{\delta\mathcal{F}}{\delta\iota(v)},
\]
with $\mathcal{F}$ a sum of four energies.
While this gradient PDE is exact, it is computationally intense.
The N-body approximation replaces the direct gradient descent
with \emph{forces} that depend only on local field values, pairwise
interactions, and sheaf-compatibility penalties.

Let each Polyxan atom be a particle with:

\[
  \text{position } x_v = \iota(v), \qquad
  \text{velocity } u_v = \dot\iota(v), \qquad
  \text{mass } m_v.
\]

We introduce an effective equation of motion of the form:
\[
  m_v \ddot{x}_v = 
    F^{(\Phi)}_v
  + F^{(v)}_v
  + F^{(S)}_v
  + F^{(\mathrm{curv})}_v
  + F^{(\mathrm{dens})}_v
  + F^{(\mathrm{sheaf})}_v.
\]

The forces correspond directly to the functional derivatives 
of $\mathcal{F}$ but expressed in a velocity/acceleration form.


% ============================================================
\section{Field-Derived Forces}
% ============================================================

\subsection{Potential Force from $\Phi$}

Define:
\[
  F^{(\Phi)}_v
    := -m_v \,\nabla_h\Phi(x_v).
\]
This places atoms inside the scalar smoothing potential:
low-$\Phi$ regions act as attractors, high-$\Phi$ regions repel.

\subsection{Directional Force from $\mathbf{v}$}

The RSVP vector field biases directional motion:
\[
  F^{(v)}_v
    := \gamma\,m_v\,\mathbf{v}(x_v).
\]
This implements the anisotropic smoothing flow described in Chapter~2.

\subsection{Entropy Force from $S$}

Entropy gradients generate a diffusion-like drift:
\[
  F^{(S)}_v
    := -\beta\,m_v\,\nabla_h(S^{-1}\nabla_h S)(x_v).
\]

This pushes atoms toward high-smoothness regions,
analogous to thermophoretic motion.


% ============================================================
\section{Hypergraph-Induced Forces}
% ============================================================

\subsection{Curvature Preservation Force}

Let $H(v)$ be the curvature estimator introduced previously.
Define the curvature force:
\[
  F^{(\mathrm{curv})}_v 
    := -\alpha_v \,\nabla H(v)^2
     = -2\alpha_v\,H(v)\,\nabla H(v).
\]

This resists bending and enforces semantic rigidity.

\subsection{Density Restoration Force}

The target density $\rho^\*$ induces a restoring force:
\[
  F^{(\mathrm{dens})}_v
    := -2\beta_v\,(\rho_\iota(v)-\rho^\*(v))
                     \nabla \rho_\iota(v).
\]

This encourages cluster formation at appropriate scales.

\subsection{Pairwise Semantic Interaction}

We may enrich the model with pairwise potentials.
For atoms $v$ and $w$:
\[
  F^{(\mathrm{pair})}_{v\leftarrow w}
  :=
  -\lambda_{vw}\,
    \nabla_x
    \Big(
      U(d(x_v,x_w))
    \Big),
\]
where $U$ is a potential (e.g.\ Lennard-Jones, Yukawa, harmonic).


% ============================================================
\section{Cohomological Force (Sheaf Alignment)}
% ============================================================

The sheaf energy was:
\[
  \mathcal{E}_{\mathrm{sheaf}}
    = \sum_{U_i} \gamma_{U_i}\,\|c^1_{i,j}\|^2
      + \gamma'_{i,j,k}\,\|c^2_{i,j,k}\|^2.
\]

The N-body approximation defines:
\[
  F^{(\mathrm{sheaf})}_v
    :=
    - \nabla_{x_v}
    \mathcal{E}_{\mathrm{sheaf}}.
\]

Explicitly, contributions arise only from neighborhoods $U_i\ni v$:
\[
  F^{(\mathrm{sheaf})}_v
    =
    -\sum_{U_i\ni v}
      \gamma_{U_i}\;
      \frac{\partial}{\partial x_v}
      \|c^1_{i,j}\|^2
    -\sum_{U_i\ni v}
      \gamma'_{i,j,k}\;
      \frac{\partial}{\partial x_v}
      \|c^2_{i,j,k}\|^2.
\]

This enforces global semantic coherence.


% ============================================================
\section{Unified N-Body System}
% ============================================================

Combining all terms:
\[
  m_v \ddot{x}_v
  =
  -m_v \nabla_h\Phi(x_v)
  + \gamma m_v\,\mathbf{v}(x_v)
  -\beta m_v\nabla_h(S^{-1}\nabla_h S)
  -2\alpha_v H(v)\,\nabla H(v)
  -2\beta_v (\rho_\iota(v)-\rho^\*)\nabla\rho_\iota(v)
  -\nabla_{x_v}\mathcal{E}_{\mathrm{sheaf}}
  + \sum_{w\neq v}F^{(\mathrm{pair})}_{v\leftarrow w}.
\]

Let:
\[
  a_v = \ddot{x}_v, \qquad
  u_v(t+\Delta t) = u_v(t) + \Delta t\,a_v, \qquad
  x_v(t+\Delta t) = x_v(t) + \Delta t\, u_v(t+\Delta t).
\]

This gives a second-order dynamical system.


% ============================================================
\section{Discrete Time Integration}
% ============================================================

A standard velocity-Verlet scheme:

\begin{align*}
  x_v(t+\tfrac{\Delta t}{2})
    &= x_v(t)
       + \tfrac{\Delta t}{2} u_v(t), \\
  u_v(t+\Delta t)
    &= u_v(t)
       + \Delta t\,a_v\big(x(t+\tfrac{\Delta t}{2})\big), \\
  x_v(t+\Delta t)
    &= x_v(t+\tfrac{\Delta t}{2})
       + \tfrac{\Delta t}{2} u_v(t+\Delta t).
\end{align*}

This retains symplecticity and minimizes drift.


% ============================================================
\section{Stability, Convergence, and Invariants}
% ============================================================

\subsection{Energy Lyapunov Property}

Define the total particle energy:
\[
  E_{\mathrm{tot}}
    = \sum_v \tfrac12 m_v\|u_v\|^2 + \mathcal{F}(\iota).
\]

N-body flow respects:
\[
  \frac{dE_{\mathrm{tot}}}{dt} \le 0
\]
when friction/dissipation are included, 
ensuring convergence to stable galaxies.


\subsection{Cluster Formation Theorem}

\begin{theorem}
If $\Phi$ has isolated minima and
$\gamma,\beta,\alpha_v,\beta_v>0$,
then trajectories satisfy:
\[
  x_v(t) \to \bigcup_k C_k,
\]
where $C_k$ are compact clusters near $\Phi$ minima.
\end{theorem}

Clusters thus correspond to stable semantic units.


\subsection{Cohomology Stabilization}

\begin{theorem}
If $\mathcal{E}_{\mathrm{sheaf}}$ dominates pairwise potentials,
then for sufficiently small $\Delta t$,
the Čech obstruction norm satisfies:
\[
  \|\mathfrak{o}(\iota_t)\| \to 0.
\]
\end{theorem}

Thus semantic inconsistency resolves over time.


% ============================================================
\section{Relation to Interpretation 3 (Cognitive Dynamics)}
% ============================================================

Under semantic interpretation:

\begin{itemize}
  \item atoms $v$ = cognitive primitives;
  \item positions $x_v$ = semantic coordinates;
  \item forces from $\Phi$ = conceptual attractors;
  \item curvature forces = conceptual rigidity or schema resistance;
  \item sheaf forces = global coherence constraints;
  \item pairwise forces = associative clustering.
\end{itemize}

Thus the N-body system is a \emph{cognitive morphogenesis engine}.


% ============================================================
\section{Summary}
% ============================================================

This chapter introduced the N-body formulation of RSVP--Polyxan simulation:

\begin{itemize}
  \item each Polyxan atom behaves as a particle influenced by RSVP fields,
        hypergraph curvature, density energies, and sheaf consistency;
  \item forces replace functional derivatives, enabling GPU acceleration;
  \item the combined system forms a stable,
        cluster-forming semantic galaxy simulator.
\end{itemize}

The next chapter will integrate visualization pipelines
for rendering RSVP fields, embedding trajectories,
and semantic galaxy structures in real time.
